\section{Cohomología de los esquemas afines}
\label{section:III.1}

\subsection{Repaso del complejo del álgebra exterior}
\label{subsection:III.1.1}

\begin{env}[1.1.1]
\label{III.1.1.1}
Sea $A$ un anillo y sea $\mathbf{f}=(f_i)_{1\leq i\leq r}$ un sistema de $r$ elementos de $A$.
El \emph{complejo del álgebra exterior $K_\bullet(\mathbf{f})$} correspondiente a $\mathbf{f}$ es un complejo de cadenas (G, I, 2.2) definido del modo siguiente: el $A$-módulo graduado $K_\bullet(\mathbf{f})$ es igual al \emph{álgebra exterior $\wedge(A^r)$}, graduada de la manera usual, y la aplicación borde es la \emph{multiplicación interior $i_\mathbf{f}$} por $\mathbf{f}$, considerado como elemento del dual $\dual{(A^r)}$; recordemos que $i_\mathbf{f}$ es una \emph{antiderivación} de grado $-1$ de $\wedge(A^r)$, y que, si $(\mathbf{e}_i)_{1\leq i\leq r}$ es la base canónica de $A^r$, se tiene $i_\mathbf{f}(\mathbf{e}_i)=f_i$; la verificación de la condición $i_\mathbf{f}\circ i_\mathbf{f}=0$ es inmediata.

Una definición equivalente es la siguiente: para cada $i$, se considera un complejo de cadenas $K_\bullet(f_i)$ definido por $K_0(f_i)=K_1(f_i)=A$, $K_n(f_i)=0$ para $n\neq 0,1$; la aplicación borde queda definida por la condición de que $d_1:A\to A$ sea la \emph{multiplicación por $f_i$}.
Se toma entonces como $K_\bullet(\mathbf{f})$ el \emph{producto tensorial $K_\bullet(f_1)\otimes K_\bullet(f_2)\otimes\cdots\otimes K_\bullet(f_r)$} (G, I, 2.7), provisto de su grado total; la verificación del isomorfismo de este complejo con el definido anteriormente es inmediata.
\end{env}

\begin{env}[1.1.2]
\label{III.1.1.2}
Para todo $A$-módulo $M$, se define el \emph{complejo de cadenas}
\[
\label{III.1.1.2.1}
  K_\bullet(\mathbf{f},M)=K_\bullet(\mathbf{f})\otimes_A M
  \tag{1.1.2.1}
\]
y el \emph{complejo de cocadenas} (G, I, 2.2)
\[
\label{III.1.1.2.2}
  K^\bullet(\mathbf{f},M)=\Hom_A(K_\bullet(\mathbf{f}),M).
  \tag{1.1.2.2}
\]

Si $g$ es una $k$-cocadena de este último complejo y se pone
\[
  g(i_1,\dots,i_k)=g(\mathbf{e}_{i_1}\wedge\cdots\wedge\mathbf{e}_{i_k}),
\]
entonces $g$ se identifica con una aplicación \emph{alternada} de $[1,r]^k$ en $M$, y de las definiciones anteriores se deduce que
\[
\label{III.1.1.2.3}
  d^k g(i_1,i_2,\dots,i_{k+1})=\sum_{h=1}^{k+1}(-1)^{h-1}f_{i_h}g(i_1,\dots,\widehat{i_h},\dots,i_{k+1}).
  \tag{1.1.2.3}
\]
\end{env}

\begin{env}[1.1.3]
\label{III.1.1.3}
\oldpage[III]{83}
De los complejos anteriores se deducen, como de costumbre, los \emph{$A$-módulos de homología y de cohomología} (G, I, 2.2)
\[
  \HH_\bullet(\mathbf{f},M)=\HH_\bullet(K_\bullet(\mathbf{f},M)),
  \tag{1.1.3.1}
\]
\[
  \HH^\bullet(\mathbf{f},M)=\HH^\bullet(K^\bullet(\mathbf{f},M)).
  \tag{1.1.3.2}
\]

Se define un \emph{$A$-isomorfismo $K_i(\mathbf{f},M)\isoto K^{r-i}(\mathbf{f},M)$} que envía cada cadena $z=\sum(\mathbf{e}_{i_1}\wedge\cdots\wedge\mathbf{e}_{i_k})\otimes z_{i_1,\dots,i_k}$ a la cocadena $g_z$ tal que $g_z(j_1,\dots,j_{r-k})=\varepsilon z_{i_1,\dots,i_k}$, donde $(j_h)_{1\leq h\leq r-k}$ es la sucesión estrictamente creciente complementaria a la sucesión estrictamente creciente $(i_h)_{1\leq h\leq k}$ en $[1,r]$ y $\varepsilon=(-1)^\nu$, siendo $\nu$ el número de inversiones de la permutación $i_1,\dots,i_k,j_1,\dots,j_{r-k}$ de $[1,r]$.
Se comprueba que $g_{dz}=d(g_z)$, lo que da un isomorfismo
\[
  \HH^i(\mathbf{f},M)\isoto\HH_{r-i}(\mathbf{f},M)\text{ para }0\leq i\leq r.
  \tag{1.1.3.3}\label{III.1.1.3.3}
\]

En este capítulo se considerarán especialmente los módulos de cohomología $\HH^\bullet(\mathbf{f},M)$.

Para un $\mathbf{f}$ dado, es inmediato (G, I, 2.1) que $M\mapsto\HH^\bullet(\mathbf{f},M)$ es un \emph{funtor cohomológico} (T, II, 2.1) de la categoría de $A$-módulos en la categoría de $A$-módulos graduados, nulo en los grados $<0$ y $>r$.
Además, se tiene
\[
  \HH^0(\mathbf{f},M)=\Hom_A(A/(\mathbf{f}),M),
  \tag{1.1.3.4}
\]
denotando por $(\mathbf{f})$ el ideal de $A$ generado por $f_1,\dots,f_r$; esto se deduce inmediatamente de (\hyperref[III.1.1.2.3]{1.1.2.3}), y es claro que $\HH^0(\mathbf{f},M)$ se identifica con el submódulo de $M$ \emph{anulado por $(\mathbf{f})$}.
De manera análoga, por (\hyperref[III.1.1.2.3]{1.1.2.3}) se tiene
\[
  \HH^r(\mathbf{f},M)=M/\bigg(\sum_{i=1}^r f_i M\bigg)=(A/(\mathbf{f}))\otimes_A M.
  \tag{1.1.3.5}\label{III.1.1.3.5}
\]

Se utilizará el siguiente resultado conocido, del que se recordará una demostración para ser completos:
\end{env}

\begin{proposition}[1.1.4]
\label{III.1.1.4}
Sean $A$ un anillo, $\mathbf{f}=(f_i)_{1\leq i\leq r}$ una familia finita de elementos de $A$ y $M$ un $A$-módulo.
Si, para $1\leq i\leq r$, la aplicación $z\mapsto f_i\cdot z$ sobre $M_{i-1}=M/(f_1 M+\cdots+f_{i-1}M)$ es inyectiva, entonces $\HH^i(\mathbf{f},M)=0$ para $i\neq r$.
\end{proposition}

\begin{proof}
Basta probar que $\HH_i(\mathbf{f},M)=0$ para todo $i>0$, en virtud de (\hyperref[III.1.1.3.3]{1.1.3.3}).
Se procede por inducción sobre $r$, siendo trivial el caso $r=0$.
Póngase $\mathbf{f}'=(f_i)_{1\leq i\leq r-1}$; esta familia satisface las condiciones del enunciado, de modo que, si se pone $L_\bullet=K_\bullet(\mathbf{f}',M)$, por hipótesis se tiene $\HH_i(L_\bullet)=0$ para $i>0$, y $\HH_0(L_\bullet)=M_{r-1}$ en virtud de (\hyperref[III.1.1.3.3]{1.1.3.3}) y (\hyperref[III.1.1.3.5]{1.1.3.5}).
Para abreviar, póngase $K_\bullet=K_\bullet(f_r)=K_0\oplus K_1$, con $K_0=K_1=A$ y $d_1:K_1\to K_0$ igual a la multiplicación por $f_r$; por definición \sref{III.1.1.1}, se tiene $K_\bullet(\mathbf{f},M)=K_\bullet\otimes_A L_\bullet$.
Se dispone del siguiente lema:

\begin{lemma}[1.1.4.1]
\label{III.1.1.4.1}
Sea $K_\bullet$ un complejo de cadenas de $A$-módulos libres, nulo salvo en las dimensiones $0$ y $1$.
Para todo complejo de cadenas $L_\bullet$ de $A$-módulos se tiene una sucesión exacta
\[
  0\to\HH_0(K_\bullet\otimes\HH_p(L_\bullet))\to\HH_p(K_\bullet\otimes L_\bullet)\to\HH_1(K_\bullet\otimes\HH_{p-1}(L_\bullet))\to 0
\]
para todo índice $p$.
\end{lemma}

\oldpage[III]{84}
Este es un caso particular de una sucesión exacta de términos de grado bajo de la sucesión espectral de K\"unneth (M, XVII, 5.2 (a) y G, I, 5.5.2); puede demostrarse directamente como sigue.
Considérense $K_0$ y $K_1$ como complejos de cadenas (nulos en las dimensiones $\neq 0$ y $\neq 1$, respectivamente); se tiene entonces una sucesión exacta de complejos
\[
  0\to K_0\otimes L_\bullet\to K_\bullet\otimes L_\bullet\to K_1\otimes L_\bullet\to 0,
\]
a la que puede aplicarse la sucesión exacta de homología
\[
  \begin{aligned}
  \cdots&\to\HH_{p+1}(K_1\otimes L_\bullet)
  \xrightarrow{\partial}\HH_p(K_0\otimes L_\bullet)
  \to\HH_p(K_\bullet\otimes L_\bullet)\\
  &\to\HH_p(K_1\otimes L_\bullet)
  \xrightarrow{\partial}\HH_{p-1}(K_0\otimes L_\bullet)\to\cdots.
  \end{aligned}
\]
Pero es evidente que $\HH_p(K_0\otimes L_\bullet)=K_0\otimes\HH_p(L_\bullet)$ y $\HH_p(K_1\otimes L_\bullet)=K_1\otimes\HH_{p-1}(L_\bullet)$ para todo $p$; además, se comprueba inmediatamente que el operador $\partial:K_1\otimes\HH_p(L_\bullet)\to K_0\otimes\HH_p(L_\bullet)$ no es otro que $d_1\otimes 1$; el lema se deduce así de la sucesión exacta anterior y de la definición de $\HH_0(K_\bullet\otimes\HH_p(L_\bullet))$ y $\HH_1(K_\bullet\otimes\HH_{p-1}(L_\bullet))$.

Establecido el lema, el final de la demostración de la Proposición \sref{III.1.1.4} es inmediato: la hipótesis de inducción y el Lema \sref{III.1.1.4.1} dan $\HH_p(K_\bullet\otimes L_\bullet)=0$ para $p\geq 2$; además, si se prueba que $\HH_1(K_\bullet,\HH_0(L_\bullet))=0$, se deduce también del Lema \sref{III.1.1.4.1} que $\HH_1(K_\bullet\otimes L_\bullet)=0$; pero, por definición, $\HH_1(K_\bullet,\HH_0(L_\bullet))$ no es otro que el núcleo de la aplicación $z\mapsto f_r\cdot z$ sobre $M_{r-1}$, y, como por hipótesis este núcleo es nulo, queda terminada la demostración.
\end{proof}

\begin{env}[1.1.5]
\label{III.1.1.5}
Sea $\mathbf{g}=(g_i)_{1\leq i\leq r}$ una segunda sucesión de $r$ elementos de $A$, y póngase $\mathbf{f}\mathbf{g}=(f_i g_i)_{1\leq i\leq r}$.
Puede definirse un homomorfismo canónico de complejos
\[
  \vphi_\mathbf{g}:K_\bullet(\mathbf{f}\mathbf{g})\to K_\bullet(\mathbf{f})
  \tag{1.1.5.1}
\]
como la extensión canónica al álgebra exterior $\wedge(A^r)$ de la aplicación $A$-lineal $(x_1,\dots,x_r)\mapsto(g_1 x_1,\dots,g_r x_r)$ de $A^r$ en sí mismo.
Para ver que se trata de un homomorfismo de complejos, basta observar, en general, que, si $u:E\to F$ es una aplicación $A$-lineal, $\mathbf{x}\in\dual{F}$ y $\mathbf{y}={}^t u(\mathbf{x})\in\dual{E}$, se tiene la fórmula
\[
  (\wedge u)\circ i_\mathbf{y}=i_\mathbf{x}\circ(\wedge u);
  \tag{1.1.5.2}
\]
en efecto, ambos miembros son antiderivaciones de $\wedge F$, y basta comprobar que coinciden sobre $F$, lo cual se deduce inmediatamente de las definiciones.

Cuando se identifica $K_\bullet(\mathbf{f})$ con el producto tensorial de los $K_\bullet(f_i)$ \sref{III.1.1.1}, $\vphi_\mathbf{g}$ es el producto tensorial de los $\vphi_{g_i}$, donde $\vphi_{g_i}$ es la identidad en grado $0$ y la multiplicación por $g_i$ en grado $1$.
\end{env}

\begin{env}[1.1.6]
\label{III.1.1.6}
En particular, para todo par de enteros $m$ y $n$ tales que $0\leq n\leq m$, se tienen homomorfismos de complejos
\[
  \vphi_{\mathbf{f}^{m-n}}:K_\bullet(\mathbf{f}^m)\to K_\bullet(\mathbf{f}^n)
  \tag{1.1.6.1}
\]
y, por consiguiente, homomorfismos
\[
  \vphi_{\mathbf{f}^{m-n}}:K^\bullet(\mathbf{f}^n,M)\to K^\bullet(\mathbf{f}^m,M),
  \tag{1.1.6.2}
\]
\[
  \vphi_{\mathbf{f}^{m-n}}:\HH^\bullet(\mathbf{f}^n,M)\to\HH^\bullet(\mathbf{f}^m,M).
  \tag{1.1.6.3}
\]

\oldpage[III]{85}
Estos últimos homomorfismos satisfacen evidentemente la condición de transitividad $\vphi_{\mathbf{f}^{m-p}}=\vphi_{\mathbf{f}^{m-n}}\circ\vphi_{\mathbf{f}^{n-p}}$ para $p\leq n\leq m$; definen, por tanto, dos \emph{sistemas inductivos} de $A$-módulos; se pone
\[
  C^\bullet((\mathbf{f}),M)=\varinjlim_n K^\bullet(\mathbf{f}^n,M),
  \tag{1.1.6.4}
\]
\[
  \HH^\bullet((\mathbf{f}),M)=\HH^\bullet(C^\bullet((\mathbf{f}),M))=\varinjlim_n\HH^\bullet(\mathbf{f}^n,M),
  \tag{1.1.6.5}
\]
deduciéndose la última igualdad del hecho de que el paso al límite inductivo conmuta con el funtor $\HH^\bullet$ (G, I, 2.1).
Se verá más adelante \sref{III.1.4.3} que $\HH^\bullet((\mathbf{f}),M)$ depende en realidad solo del \emph{ideal $(\mathbf{f})$} de $A$ (e incluso solo de la topología $(\mathbf{f})$-preádica sobre $A$), lo que justifica las notaciones.

Es claro que $M\mapsto C^\bullet((\mathbf{f}),M)$ es un funtor exacto $A$-lineal y que $M\mapsto\HH^\bullet((\mathbf{f}),M)$ es un funtor cohomológico.
\end{env}

\begin{env}[1.1.7]
\label{III.1.1.7}
Pónganse $\mathbf{f}=(f_i)\in A^r$ y $\mathbf{g}=(g_i)\in A^r$; denótese por $e_\mathbf{g}$ la multiplicación por la izquierda por el vector $\mathbf{g}\in A^r$ sobre el álgebra exterior $\wedge(A^r)$; se sabe que se tiene la \emph{fórmula de homotopía}
\[
  i_\mathbf{f}e_\mathbf{g}+e_\mathbf{g}i_\mathbf{f}=\langle\mathbf{g},\mathbf{f}\rangle 1
  \tag{1.1.7.1}
\]
en el $A$-módulo $\wedge(A^r)$ ($1$ denota el automorfismo identidad de $\wedge(A^r)$); esta relación implica también que \emph{en el complejo $K_\bullet(\mathbf{f})$} se tiene
\[
  de_\mathbf{g}+e_\mathbf{g}d=\langle\mathbf{g},\mathbf{f}\rangle 1.
  \tag{1.1.7.2}
\]

Si el ideal $(\mathbf{f})$ es igual a $A$, existe $\mathbf{g}\in A^r$ tal que $\langle\mathbf{g},\mathbf{f}\rangle=\sum_{i=1}^r g_i f_i=1$.
Por consiguiente (G, I, 2.4):
\end{env}

\begin{proposition}[1.1.8]
\label{III.1.1.8}
Supóngase que el ideal $(\mathbf{f})$ generado por los $f_i$ es igual a $A$.
Entonces el complejo $K_\bullet(\mathbf{f})$ es homotópicamente trivial, al igual que los complejos $K_\bullet(\mathbf{f},M)$ y $K^\bullet(\mathbf{f},M)$ para todo $A$-módulo $M$.
\end{proposition}

\begin{corollary}[1.1.9]
\label{III.1.1.9}
Si $(\mathbf{f})=A$, entonces $\HH^\bullet(\mathbf{f},M)=0$ y $\HH^\bullet((\mathbf{f}),M)=0$ para todo $A$-módulo $M$.
\end{corollary}

\begin{proof}
En efecto, se tiene entonces $(\mathbf{f}^n)=A$ para todo $n$.
\end{proof}

\begin{remark}[1.1.10]
\label{III.1.1.10}
Con las mismas notaciones anteriores, póngase $X=\Spec(A)$ y sea $Y$ el subpreesquema cerrado de $X$ definido por el ideal $(\mathbf{f})$.
Se probará en \textsection9 que $\HH^\bullet((\mathbf{f}),M)$ es isomorfo a la cohomología $\HH_Y^\bullet(X,\widetilde{M})$ correspondiente al antifiltro $\Phi$ de subconjuntos cerrados de $Y$ (T, 3.2).
Se mostrará también que la Proposición \sref{III.1.2.3}, aplicada a $X$ y a $\sh{F}=\widetilde{M}$, es un caso particular de una sucesión exacta de cohomología
\[
  \cdots\to\HH_Y^p(X,\sh{F})\to\HH^p(X,\sh{F})\to\HH^p(X\setmin Y,\sh{F})\to\HH_Y^{p+1}(X,\sh{F})\to\cdots.
\]
\end{remark}

\subsection{Cohomología de \v Cech de un recubrimiento abierto}
\label{subsection:III.1.2}

\begin{notation}[1.2.1]
\label{III.1.2.1}
En esta sección se denota por:
\begin{enumerate}
  \item $X$ un preesquema;
  \item $\sh{F}$ un $\sh{O}_X$-módulo cuasi-coherente;
\oldpage[III]{86}
  \item $A=\Gamma(X,\sh{O}_X)$, $M=\Gamma(X,\sh{F})$;
  \item $\mathbf{f}=(f_i)_{1\leq i\leq r}$ un sistema finito de elementos de $A$;
  \item $U_i=X_{f_i}$, el abierto \sref[0]{0.5.5.2} de los $x\in X$ tales que $f_i(x)\neq 0$;
  \item $U=\bigcup_{i=1}^r U_i$;
  \item $\mathfrak{U}$ el recubrimiento $(U_i)_{1\leq i\leq r}$ de $U$.
\end{enumerate}
\end{notation}

\begin{env}[1.2.2]
\label{III.1.2.2}
Supóngase que $X$ es un preesquema cuyo espacio subyacente es \emph{noetheriano}, o bien un \emph{esquema} cuyo espacio subyacente es \emph{cuasicompacto}.
Se sabe entonces \sref[I]{I.9.3.3} que $\Gamma(U_i,\sh{F})=M_{f_i}$.
Póngase
\[
  U_{i_0 i_1\cdots i_p}=\bigcap_{k=0}^p U_{i_k}=X_{f_{i_0}f_{i_1}\cdots f_{i_p}}
\]
\sref[0]{0.5.5.3}; se tiene, por tanto, también
\[
  \Gamma(U_{i_0 i_1\cdots i_p},\sh{F})=M_{f_{i_0}f_{i_1}\cdots f_{i_p}}.
  \tag{1.2.2.1}
\]

Por \sref[0]{0.1.6.1}, $M_{f_{i_0}f_{i_1}\cdots f_{i_p}}$ se identifica con el límite inductivo $\varinjlim_n M_{i_0 i_1\cdots i_p}^{(n)}$, donde el sistema inductivo está formado por los $M_{i_0 i_1\cdots i_p}^{(n)}=M$, siendo los homomorfismos $\vphi_{nm}:M_{i_0 i_1\cdots i_p}^{(m)}\to M_{i_0 i_1\cdots i_p}^{(n)}$ la multiplicación por $(f_{i_0}f_{i_1}\cdots f_{i_p})^{n-m}$ para $m\leq n$.
Se denota por $C_n^p(M)$ el conjunto de aplicaciones \emph{alternadas} de $[1,r]^{p+1}$ en $M$ (para todo $n$); estos $A$-módulos forman también un sistema inductivo respecto de los $\vphi_{nm}$.
Si $C^p(\mathfrak{U},\sh{F})$ es el grupo de las $p$-cocadenas \emph{alternadas} de \v Cech relativas al recubrimiento $\mathfrak{U}$, con coeficientes en $\sh{F}$ (G, II, 5.1), de lo anterior se deduce que puede escribirse
\[
  C^p(\mathfrak{U},\sh{F})=\varinjlim_n C_n^p(M).
  \tag{1.2.2.2}
\]

Con las notaciones de \sref{III.1.1.2}, $C_n^p(M)$ se identifica con $K^{p+1}(\mathbf{f}^n,M)$, y la aplicación $\vphi_{nm}$ se identifica con la aplicación $\vphi_{\mathbf{f}^{n-m}}$ definida en \sref{III.1.1.6}.
Se tiene así, para todo $p\geq 0$, un isomorfismo canónico funtorial
\[
  C^p(\mathfrak{U},\sh{F})\isoto C^{p+1}((\mathbf{f}),M).
  \tag{1.2.2.3}\label{III.1.2.2.3}
\]

Además, la fórmula (\hyperref[III.1.1.2.3]{1.1.2.3}) y la definición de la cohomología de un recubrimiento (G, II, 5.1) muestran que los isomorfismos (\hyperref[III.1.2.2.3]{1.2.2.3}) son compatibles con las aplicaciones coborde.
\end{env}

\begin{proposition}[1.2.3]
\label{III.1.2.3}
Si $X$ es un preesquema cuyo espacio subyacente es noetheriano, o un esquema cuyo espacio subyacente es cuasicompacto, existe un isomorfismo canónico funtorial en $\sh{F}$
\[
  \HH^p(\mathfrak{U},\sh{F})\isoto\HH^{p+1}((\mathbf{f}),M)\text{ para }p\geq 1.
  \tag{1.2.3.1}\label{III.1.2.3.1}
\]

Además, se tiene una sucesión exacta funtorial en $\sh{F}$
\[
  0\to\HH^0((\mathbf{f}),M)\to M\to\HH^0(\mathfrak{U},\sh{F})\to\HH^1((\mathbf{f}),M)\to 0.
  \tag{1.2.3.2}\label{III.1.2.3.2}
\]
\end{proposition}

\oldpage[III]{87}
\begin{proof}
Los isomorfismos (\hyperref[III.1.2.3.1]{1.2.3.1}) son consecuencias inmediatas de lo visto en \sref{III.1.2.2}.
Por otra parte, se tiene $C^0(\mathfrak{U},\sh{F})=C^1((\mathbf{f}),M)$; por consiguiente, $\HH^0(\mathfrak{U},\sh{F})$ se identifica con el subgrupo de los $1$-cociclos de $C^1((\mathbf{f}),M)$; como $M=C^0((\mathbf{f}),M)$, la sucesión exacta (\hyperref[III.1.2.3.2]{1.2.3.2}) no es otra que la dada por la definición de los grupos de cohomología $\HH^0((\mathbf{f}),M)$ y $\HH^1((\mathbf{f}),M)$.
\end{proof}

\begin{corollary}[1.2.4]
\label{III.1.2.4}
Supóngase que los $X_{f_i}$ son cuasicompactos y que existen $g_i\in\Gamma(U,\sh{O}_X)$ tales que $\sum_i g_i(f_i|U)=1|U$.
Entonces, para todo $(\sh{O}_X|U)$-módulo cuasi-coherente $\sh{G}$, se tiene $\HH^p(\mathfrak{U},\sh{G})=0$ para $p>0$; si además $U=X$, el homomorfismo canónico (\hyperref[III.1.2.3.2]{1.2.3.2}) $\Gamma(X,\sh{G})\to\HH^0(\mathfrak{U},\sh{G})$ es biyectivo.
\end{corollary}

\begin{proof}
Como, por hipótesis, los $U_i=X_{f_i}$ son cuasicompactos, también lo es $U$, y puede reducirse al caso $U=X$; poniendo $N=\Gamma(X,\sh{G})$, la hipótesis implica entonces que $\HH^p((\mathbf{f}),N)=0$ para todo $p\geq 0$ \sref{III.1.1.9}.
El corolario se deduce inmediatamente de (\hyperref[III.1.2.3.1]{1.2.3.1}) y (\hyperref[III.1.2.3.2]{1.2.3.2}).
\end{proof}

Obsérvese que, puesto que $\HH^0(\mathfrak{U},\sh{F})=\HH^0(U,\sh{F})$ (G, II, 5.2.2), se ha demostrado de nuevo \sref[I]{I.1.3.7} como caso particular.

\begin{remark}[1.2.5]
\label{III.1.2.5}
Supóngase que $X$ es un \emph{esquema afín}; entonces los $U_i=X_{f_i}=D(f_i)$ son abiertos afines, al igual que los $U_{i_0 i_1\cdots i_p}$ (pero $U$ no es necesariamente afín).
En este caso, los funtores $\Gamma(X,\sh{F})$ y $\Gamma(U_{i_0 i_1\cdots i_p},\sh{F})$ son exactos en $\sh{F}$ \sref[I]{I.1.3.11}.
Si se tiene una sucesión exacta $0\to\sh{F}'\to\sh{F}\to\sh{F}''\to 0$ de $\sh{O}_X$-módulos cuasi-coherentes, la sucesión de complejos
\[
  0\to C^\bullet(\mathfrak{U},\sh{F}')\to C^\bullet(\mathfrak{U},\sh{F})\to C^\bullet(\mathfrak{U},\sh{F}'')\to 0
\]
es exacta y da, por tanto, una sucesión exacta de cohomología
\[
  \cdots\to\HH^p(\mathfrak{U},\sh{F}')\to\HH^p(\mathfrak{U},\sh{F})\to\HH^p(\mathfrak{U},\sh{F}'')\xrightarrow{\partial}\HH^{p+1}(\mathfrak{U},\sh{F}')\to\cdots.
\]
Por otra parte, si se pone $M'=\Gamma(X,\sh{F}')$ y $M''=\Gamma(X,\sh{F}'')$, la sucesión $0\to M'\to M\to M''\to 0$ es exacta; como $C^\bullet((\mathbf{f}),M)$ es un funtor exacto en $M$, se tiene asimismo la sucesión exacta de cohomología
\[
  \cdots\to\HH^p((\mathbf{f}),M')\to\HH^p((\mathbf{f}),M)\to\HH^p((\mathbf{f}),M'')\xrightarrow{\partial}\HH^{p+1}((\mathbf{f}),M')\to\cdots.
\]

Así pues, como el diagrama
\[
  \xymatrix{
    0\ar[r] &
    C^\bullet(\mathfrak{U},\sh{F}')\ar[r]\ar[d] &
    C^\bullet(\mathfrak{U},\sh{F})\ar[r]\ar[d] &
    C^\bullet(\mathfrak{U},\sh{F}'')\ar[r]\ar[d] &
    0\\
    0\ar[r] &
    C^\bullet((\mathbf{f}),M')\ar[r] &
    C^\bullet((\mathbf{f}),M)\ar[r] &
    C^\bullet((\mathbf{f}),M'')\ar[r] &
    0
  }
\]
\oldpage[III]{88}
es conmutativo, se concluye que los diagramas
\[
  \xymatrix{
    \HH^p(\mathfrak{U},\sh{F}'')\ar[r]^\partial\ar[d] &
    \HH^{p+1}(\mathfrak{U},\sh{F}')\ar[d]\\
    \HH^{p+1}((\mathbf{f}),M'')\ar[r]^\partial &
    \HH^{p+2}((\mathbf{f}),M')
  }
  \tag{1.2.5.1}
\]
son conmutativos para todo $p$~(G, I, 2.1.1).
\end{remark}

\subsection{Cohomología de un esquema afín}
\label{subsection:III.1.3}

\begin{theorem}[1.3.1]
\label{III.1.3.1}
Sea $X$ un esquema afín.
Para todo $\sh{O}_X$-módulo cuasi-coherente $\sh{F}$, se tiene $\HH^p(X,\sh{F})=0$ para todo $p>0$.
\end{theorem}

\begin{proof}
Sea $\mathfrak{U}$ un recubrimiento finito de $X$ por los abiertos afines $X_{f_i}=D(f_i)$ ($1\leq i\leq r$); se sabe entonces que el ideal de $A=\Gamma(X,\sh{O}_X)$ generado por los $f_i$ es igual a $A$.
Del Corolario~\sref{III.1.2.4} se concluye, pues, que $\HH^p(\mathfrak{U},\sh{F})=0$ para $p>0$.
Como existen recubrimientos finitos de $X$ por abiertos afines arbitrariamente finos \sref[I]{I.1.1.10}, la definición de la cohomología de \v Cech (G, II, 5.8) muestra que también $\CHH^p(X,\sh{F})=0$ para $p>0$.
Pero esto se aplica asimismo a todo preesquema $X_f$, para $f\in A$ \sref[I]{I.1.3.6}; por tanto, $\CHH^p(X_f,\sh{F})=0$ para $p>0$.
Como $X_f\cap X_g=X_{fg}$, se deduce también que $\HH^p(X,\sh{F})=0$ para todo $p>0$, en virtud de (G, II, 5.9.2).
\end{proof}

\begin{corollary}[1.3.2]
\label{III.1.3.2}
Sea $Y$ un preesquema y sea $f:X\to Y$ un morfismo afín \sref[II]{II.1.6.1}.
Para todo $\sh{O}_X$-módulo cuasi-coherente $\sh{F}$, se tiene $\RR^q f_*(\sh{F})=0$ para $q>0$.
\end{corollary}

\begin{proof}
Por definición, $\RR^q f_*(\sh{F})$ es el $\sh{O}_Y$-módulo asociado al prehaz $U\mapsto\HH^q(f^{-1}(U),\sh{F})$, donde $U$ recorre los abiertos de $Y$.
Pero los abiertos afines forman una base de $Y$, y, para un abierto tal $U$, $f^{-1}(U)$ es afín \sref[II]{II.1.3.2}; por tanto, $\HH^q(f^{-1}(U),\sh{F})=0$ por el Teorema~\sref{III.1.3.1}, lo que prueba el corolario.
\end{proof}

\begin{corollary}[1.3.3]
\label{III.1.3.3}
Sea $Y$ un preesquema y sea $f:X\to Y$ un morfismo afín.
Para todo $\sh{O}_X$-módulo cuasi-coherente $\sh{F}$, el homomorfismo canónico $\HH^p(Y,f_*(\sh{F}))\to\HH^p(X,\sh{F})$ (\textbf{0},~12.1.3.1) es biyectivo para todo $p$.
\end{corollary}

\begin{proof}
Basta (por \sref[0]{0.12.1.7}) mostrar que los homomorfismos de borde $''E_2^{p0}=\HH^p(Y,f_*(\sh{F}))\to\HH^p(X,\sh{F})$ de la segunda sucesión espectral del funtor compuesto $\Gamma f_*$ son biyectivos.
Pero el término $E_2$ de esta sucesión espectral viene dado por $''E_2^{pq}=\HH^p(Y,\RR^q f_*(\sh{F}))$ (G, II, 4.17.1); por el Corolario~\sref{III.1.3.2}, $''E_2^{pq}=0$ para $q>0$, y la sucesión espectral degenera; de ahí la afirmación \sref[0]{0.11.1.6}.
\end{proof}

\begin{corollary}[1.3.4]
\label{III.1.3.4}
Sean $f:X\to Y$ un morfismo afín y $g:Y\to Z$ un morfismo.
\oldpage[III]{89}
Para todo $\sh{O}_X$-módulo cuasi-coherente $\sh{F}$, el homomorfismo canónico $\RR^p g_*(f_*(\sh{F}))\to\RR^p(g\circ f)_*(\sh{F})$ (\textbf{0},~12.2.5.1) es biyectivo para todo $p$.
\end{corollary}

\begin{proof}
Basta observar que, según el Corolario~\sref{III.1.3.3}, para todo abierto afín $W$ de $Z$, el homomorfismo canónico $\HH^p(g^{-1}(W),f_*(\sh{F}))\to\HH^p(f^{-1}(g^{-1}(W)),\sh{F})$ es biyectivo; esto prueba que el homomorfismo de prehaces que define el homomorfismo canónico $\RR^p g_*(f_*(\sh{F}))\to\RR^p(g\circ f)_*(\sh{F})$ es biyectivo \sref[0]{0.12.2.5}.
\end{proof}

\subsection{Aplicación a la cohomología de preesquemas arbitrarios}
\label{subsection:III.1.4}

\begin{proposition}[1.4.1]
\label{III.1.4.1}
Sea $X$ un esquema y sea $\mathfrak{U}=(U_\alpha)$ un recubrimiento de $X$ por abiertos afines.
Para todo $\sh{O}_X$-módulo cuasi-coherente $\sh{F}$, los módulos de cohomología $\HH^\bullet(X,\sh{F})$ y $\HH^\bullet(\mathfrak{U},\sh{F})$ (sobre $\Gamma(X,\sh{O}_X)$) son canónicamente isomorfos.
\end{proposition}

\begin{proof}
Como $X$ es un esquema, toda intersección finita $V$ de abiertos del recubrimiento $\mathfrak{U}$ es afín \sref[I]{I.5.5.6}; por tanto, $\HH^q(V,\sh{F})=0$ para $q\geq 1$ por el Teorema~\sref{III.1.3.1}.
La proposición se deduce entonces de un teorema de Leray (G, II, 5.4.1).
\end{proof}

\begin{remark}[1.4.2]
\label{III.1.4.2}
Obsérvese que el resultado de la Proposición~\sref{III.1.4.1} sigue siendo válido cuando las intersecciones finitas de los conjuntos $U_\alpha$ son afines, incluso si no se supone necesariamente que $X$ sea un esquema.
\end{remark}

\begin{corollary}[1.4.3]
\label{III.1.4.3}
Sea $X$ un esquema de espacio subyacente cuasicompacto, sea $A=\Gamma(X,\sh{O}_X)$ y sea $\mathbf{f}=(f_i)_{1\leq i\leq r}$ una sucesión finita de elementos de $A$ tal que los $X_{f_i}$ (notación de \sref{III.1.2.1}) sean afines.
Entonces (con las notaciones de \sref{III.1.2.1}), para todo $\sh{O}_X$-módulo cuasi-coherente $\sh{F}$, se tiene un isomorfismo canónico funtorial en $\sh{F}$
\[
  \HH^q(U,\sh{F})\isoto\HH^{q+1}((\mathbf{f}),M)\text{ para }q\geq 1,
  \tag{1.4.3.1}\label{III.1.4.3.1}
\]
y una sucesión exacta funtorial en $\sh{F}$
\[
  0\to\HH^0((\mathbf{f}),M)\to M\to\HH^0(U,\sh{F})\to\HH^1((\mathbf{f}),M)\to 0.
  \tag{1.4.3.2}
\]
\end{corollary}

\begin{proof}
Esto se deduce inmediatamente de las Proposiciones~\sref{III.1.4.1} y \sref{III.1.2.3}.
\end{proof}

\begin{env}[1.4.4]
\label{III.1.4.4}
Si $X$ es un \emph{esquema afín}, de la Observación~\sref{III.1.2.5} y de la Proposición~\sref{III.1.4.1} se deduce que, para todo $q\geq 0$, los diagramas
\[
  \xymatrix{
    \HH^q(U,\sh{F}'')\ar[r]^\partial\ar[d] &
    \HH^{q+1}(U,\sh{F}')\ar[d]\\
    \HH^{q+1}((\mathbf{f}),M'')\ar[r]^\partial &
    \HH^{q+2}((\mathbf{f}),M')
  }
  \tag{1.4.4.1}
\]
correspondientes a una sucesión exacta $0\to\sh{F}'\to\sh{F}\to\sh{F}''\to 0$ de $\sh{O}_X$-módulos cuasi-coherentes (con las notaciones de la Observación~\sref{III.1.2.5}) son conmutativos.
\end{env}

\begin{proposition}[1.4.5]
\label{III.1.4.5}
\oldpage[III]{90}
Sea $X$ un esquema cuasicompacto, sea $\sh{L}$ un $\sh{O}_X$-módulo invertible, y considérese el anillo graduado $A_\bullet=\Gamma_\bullet(\sh{L})$ \sref[0]{0.5.4.6}; entonces $\HH^\bullet(\sh{F},\sh{L})=\bigoplus_{n\in\bb{Z}}\HH^\bullet(X,\sh{F}\otimes\sh{L}^{\otimes n})$ es un $A_\bullet$-módulo graduado y, para todo $f\in A_n$, se tiene un isomorfismo canónico
\[
  \HH^\bullet(X_f,\sh{F})\isoto(\HH^\bullet(\sh{F},\sh{L}))_{(f)}
  \tag{1.4.5.1}\label{III.1.4.5.1}
\]
de $(A_\bullet)_{(f)}$-módulos.
\end{proposition}

\begin{proof}
Como $X$ es un esquema cuasicompacto, puede calcularse la cohomología de todos los $\sh{O}_X$-módulos $\sh{F}\otimes\sh{L}^{\otimes n}$ utilizando un mismo recubrimiento finito $\mathfrak{U}=(U_i)$, formado por abiertos afines tales que la restricción $\sh{L}|U_i$ sea isomorfa a $\sh{O}_X|U_i$ para cada $i$ \sref{III.1.4.1}.
Es entonces inmediato que los $U_i\cap X_f$ son abiertos afines \sref[I]{I.1.3.6}, y puede calcularse así la cohomología $\HH^\bullet(X_f,\sh{F}\otimes\sh{L}^{\otimes n})$ utilizando el recubrimiento $\mathfrak{U}|X_f=(U_i\cap X_f)$ \sref{III.1.4.1}.
Es inmediato que, para todo $f\in A_n$, la multiplicación por $f$ define un homomorfismo $C^\bullet(\mathfrak{U},\sh{F}\otimes\sh{L}^m)\to C^\bullet(\mathfrak{U},\sh{F}\otimes\sh{L}^{\otimes(m+n)})$ y, por tanto, un homomorfismo $\HH^\bullet(\mathfrak{U},\sh{F}\otimes\sh{L}^{\otimes m})\to\HH^\bullet(\mathfrak{U},\sh{F}\otimes\sh{L}^{\otimes(m+n)})$, lo que establece la primera afirmación.
Por otra parte, para un $f\in A_n$ dado, de \sref[I]{I.9.3.2} se deduce que se tiene un isomorfismo de complejos de $(A_\bullet)_{(f)}$-módulos
\[
  C^\bullet(\mathfrak{U}|X_f,\sh{F})\isoto\bigg(\!\!C^\bullet\bigg(\mathfrak{U},\bigoplus_{n\in\bb{Z}}\sh{F}\otimes\sh{L}^{\otimes n}\bigg)\!\!\bigg)_{(f)},
\]
teniendo en cuenta \sref[I]{I.1.3.9}[ii].
Al pasar a la cohomología de estos complejos se obtiene el isomorfismo (\hyperref[III.1.4.5.1]{1.4.5.1}), recordando que el funtor $M\mapsto M_{(f)}$ es exacto sobre la categoría de $A_\bullet$-módulos graduados.
\end{proof}

\begin{corollary}[1.4.6]
\label{III.1.4.6}
Supóngase que se satisfacen las hipótesis de la Proposición~\sref{III.1.4.5} y, además, que $\sh{L}=\sh{O}_X$.
Si se pone $A=\Gamma(X,\sh{O}_X)$, entonces, para todo $f\in A$, se tiene un isomorfismo canónico $\HH^\bullet(X_f,\sh{F})\isoto(\HH^\bullet(X,\sh{F}))_f$ de $A_f$-módulos.
\end{corollary}

\begin{corollary}[1.4.7]
\label{III.1.4.7}
Sea $X$ un esquema cuasicompacto y sea $f$ un elemento de $\Gamma(X,\sh{O}_X)$.
\begin{enumerate}
  \item[{\rm(i)}] Supóngase que el abierto $X_f$ es afín.
    Entonces, para todo $\sh{O}_X$-módulo cuasi-coherente $\sh{F}$, todo $i>0$ y todo $\xi\in\HH^i(X,\sh{F})$, existe un entero $n>0$ tal que $f^n\xi=0$.
  \item[{\rm(ii)}] Recíprocamente, supóngase que $X_f$ es cuasicompacto y que, para todo haz cuasi-coherente de ideales $\sh{J}$ de $\sh{O}_X$ y todo $\zeta\in\HH^1(X,\sh{J})$, existe un $n>0$ tal que $f^n\zeta=0$.
    Entonces $X_f$ es afín.
\end{enumerate}
\end{corollary}

\begin{proof}
\medskip\noindent
\begin{enumerate}
  \item[(i)] Si $X_f$ es afín, se tiene $\HH^i(X_f,\sh{F})=0$ para todo $i>0$ \sref{III.1.3.1}, de modo que la afirmación se deduce directamente del Corolario~\sref{III.1.4.6}.
  \item[(ii)] En virtud del criterio de Serre \sref[II]{II.5.2.1}, basta probar que, para todo haz cuasi-coherente de ideales $\sh{K}$ de $\sh{O}_X|X_f$, se tiene $\HH^1(X_f,\sh{K})=0$.
    Como $X_f$ es un abierto cuasicompacto de un esquema cuasicompacto $X$, existe un haz cuasi-coherente de ideales $\sh{J}$ de $\sh{O}_X$ tal que $\sh{K}=\sh{J}|X_f$ \sref[I]{I.9.4.2}.
    Según el Corolario~\sref{III.1.4.6}, se tiene $\HH^1(X_f,\sh{K})=(\HH^1(X,\sh{J}))_f$, y la hipótesis implica que el miembro de la derecha es nulo, lo que da la afirmación.
\end{enumerate}
\end{proof}

\begin{remark}[1.4.8]
\label{III.1.4.8}
Obsérvese que el Corolario~\sref{III.1.4.7}[i] da una demostración más sencilla de la relación (\textbf{II},~4.5.13.2).
\end{remark}

\begin{lemma}[1.4.9]
\label{III.1.4.9}
\oldpage[III]{91}
Sea $X$ un esquema cuasicompacto, sea $\mathfrak{U}=(U_i)_{1\leq i\leq n}$ un recubrimiento finito de $X$ por abiertos afines y sea $\sh{F}$ un $\sh{O}_X$-módulo cuasi-coherente.
El complejo de haces $\sh{C}^\bullet(\mathfrak{U},\sh{F})$ definido por el recubrimiento $\mathfrak{U}$ (G, II, 5.2) es entonces un $\sh{O}_X$-módulo cuasi-coherente.
\end{lemma}

\begin{proof}
De las definiciones (G, II, 5.2) se deduce que $\sh{C}^p(\mathfrak{U},\sh{F})$ es la suma directa de los haces imagen directa de los $\sh{F}|U_{i_0\cdots i_p}$ por la inyección canónica $U_{i_0\cdots i_p}\to X$.
La hipótesis de que $X$ es un esquema implica que estas inyecciones son morfismos afines \sref[I]{I.5.5.6}; por tanto, los $\sh{C}^p(\mathfrak{U},\sh{F})$ son cuasi-coherentes \sref[II]{II.1.2.6}.
\end{proof}

\begin{proposition}[1.4.10]
\label{III.1.4.10}
Sea $u:X\to Y$ un morfismo separado y cuasicompacto.
Para todo $\sh{O}_X$-módulo cuasi-coherente $\sh{F}$, los $\RR^q u_*(\sh{F})$ son $\sh{O}_Y$-módulos cuasi-coherentes.
\end{proposition}

\begin{proof}
La cuestión es local sobre $Y$, de modo que puede suponerse que $Y$ es afín.
Entonces $X$ es una unión finita de abiertos afines $U_i$ ($1\leq i\leq n$); sea $\mathfrak{U}$ el recubrimiento $(U_i)$.
Además, como $Y$ es un esquema, de \sref[I]{I.5.5.10} se deduce que, para todo abierto afín $V\subset Y$, la inyección canónica $u^{-1}(V)\to X$ es un morfismo afín; se concluye (Proposición~\sref{III.1.4.1} y (G, II, 5.2)) que se tiene un isomorfismo canónico
\[
  \HH^\bullet(u^{-1}(V),\sh{F})\isoto\HH^\bullet(\Gamma(V,\sh{K}^\bullet)),
  \tag{1.4.10.1}\label{III.1.4.10.1}
\]
donde se pone $\sh{K}^\bullet=u_*(\sh{C}^\bullet(\mathfrak{U},\sh{F}))$.
Según el Lema~\sref{III.1.4.9} y \sref[I]{I.9.2.2}, $\sh{K}^\bullet$ es un $\sh{O}_Y$-módulo cuasi-coherente; además, constituye un \emph{complejo de haces}, puesto que también lo es $\sh{C}^\bullet(\mathfrak{U},\sh{F})$.
De la definición de la cohomología $\sh{H}^\bullet(\sh{K}^\bullet)$ (G, II, 4.1) se deduce entonces que esta última está formada por $\sh{O}_Y$-módulos cuasi-coherentes \sref[I]{I.4.1.1}.
Como (para $V$ afín en $Y$) el funtor $\Gamma(V,\sh{G})$ es exacto en $\sh{G}$ sobre la categoría de $\sh{O}_Y$-módulos cuasi-coherentes, se tiene (G, II, 4.1)
\[
  \HH^\bullet(\Gamma(V,\sh{K}^\bullet))=\Gamma(V,\sh{H}^\bullet(\sh{K}^\bullet)).
  \tag{1.4.10.2}\label{III.1.4.10.2}
\]

Finalmente, obsérvese que de la definición del homomorfismo canónico
\[
  \HH^\bullet(\mathfrak{U},\sh{F})\to\HH^\bullet(X,\sh{F}),
\]
dado en (G, II, 5.2), se deduce que, si $V'\subset V$ es otro abierto afín de $Y$, el
diagrama
\[
  \xymatrix{
    \HH^\bullet(u^{-1}(V),\sh{F})\ar[r]^\sim\ar[d] &
    \HH^\bullet(\Gamma(V,\sh{K}^\bullet))\ar[d]\\
    \HH^\bullet(u^{-1}(V'),\sh{F})\ar[r]^\sim &
    \HH^\bullet(\Gamma(V',\sh{K}^\bullet))
  }
\]
es conmutativo.
De lo anterior se concluye que los isomorfismos (\hyperref[III.1.4.10.1]{1.4.10.1}) definen un isomorfismo de $\sh{O}_Y$-módulos
\[
  \RR^\bullet u_*(\sh{F})\isoto\sh{H}^\bullet(\sh{K}^\bullet),
  \tag{1.4.10.3}\label{III.1.4.10.3}
\]
\oldpage[III]{92}
y, por consiguiente, $\RR^\bullet u_*(\sh{F})$ es cuasi-coherente.
\end{proof}

Además, de (\hyperref[III.1.4.10.3]{1.4.10.3}), (\hyperref[III.1.4.10.2]{1.4.10.2}) y (\hyperref[III.1.4.10.1]{1.4.10.1}) se deduce que:
\begin{corollary}[1.4.11]
\label{III.1.4.11}
Bajo las hipótesis de la Proposición~\sref{III.1.4.10}, para todo abierto afín $V$ de $Y$, el homomorfismo canónico
\[
  \HH^q(u^{-1}(V),\sh{F})\to\Gamma(V,\RR^q u_*(\sh{F}))
  \tag{1.4.11.1}\label{III.1.4.11.1}
\]
es un isomorfismo para todo $q\geq 0$.
\end{corollary}

\begin{corollary}[1.4.12]
\label{III.1.4.12}
Supóngase que se satisfacen las hipótesis de la Proposición~\sref{III.1.4.10} y, además, que $Y$ es cuasicompacto.
Entonces existe un entero $r>0$ tal que, para todo $\sh{O}_X$-módulo cuasi-coherente $\sh{F}$ y todo entero $q>r$, se tiene $\RR^q u_*(\sh{F})=0$.
Si $Y$ es afín, puede tomarse como $r$ un entero tal que exista un recubrimiento de $X$ formado por $r$ abiertos afines.
\end{corollary}

\begin{proof}
Como puede recubrirse $Y$ por un número finito de abiertos afines, en virtud del Corolario~\sref{III.1.4.11} puede reducirse la demostración a la segunda afirmación.
Si $\mathfrak{U}$ es un recubrimiento de $X$ por $r$ abiertos afines, se tiene $\HH^q(\mathfrak{U},\sh{F})=0$ para $q>r$, puesto que las cocadenas de $C^q(\mathfrak{U},\sh{F})$ son alternadas; la afirmación se deduce entonces de la Proposición~\sref{III.1.4.1}.
\end{proof}

\begin{corollary}[1.4.13]
\label{III.1.4.13}
Supóngase que se satisfacen las hipótesis de la Proposición~\sref{III.1.4.10} y, además, que $Y=\Spec(A)$ es afín.
Entonces, para todo $\sh{O}_X$-módulo cuasi-coherente $\sh{F}$ y todo $f\in A$, se tiene
\[
  \Gamma(Y_f,\RR^q u_*(\sh{F}))=(\Gamma(Y,\RR^q u_*(\sh{F})))_f
\]
salvo isomorfismo canónico.
\end{corollary}

\begin{proof}
Esto se deduce del hecho de que $\RR^q u_*(\sh{F})$ es un $\sh{O}_Y$-módulo cuasi-coherente \sref[I]{I.1.3.7}.
\end{proof}

\begin{proposition}[1.4.14]
\label{III.1.4.14}
Sean $f:X\to Y$ un morfismo separado y cuasicompacto y $g:Y\to Z$ un morfismo afín.
Para todo $\sh{O}_X$-módulo cuasi-coherente $\sh{F}$, el homomorfismo canónico $\RR^p(g\circ f)_*(\sh{F})\to g_*(\RR^p f_*(\sh{F}))$ (\textbf{0},~12.2.5.2) es biyectivo para todo $p$.
\end{proposition}

\begin{proof}
Para todo abierto afín $W$ de $Z$, $g^{-1}(W)$ es un abierto afín de $Y$.
El homomorfismo de prehaces que define el homomorfismo canónico
\[
  \RR^p(g\circ f)_*(\sh{F})\to g_*(\RR^p f_*(\sh{F}))
\]
\sref[0]{0.12.2.5} es, por tanto, biyectivo en virtud del Corolario~\sref{III.1.4.11}.
\end{proof}

\begin{proposition}[1.4.15]
\label{III.1.4.15}
Sean $u:X\to Y$ un morfismo separado cuasicompacto y $v:Y'\to Y$ un morfismo plano de preesquemas \sref[0]{0.6.7.1}; sea $u'=u_{(Y')}$, de modo que se tiene el diagrama conmutativo
\[
  \xymatrix{
    X\ar[d]_u &
    X'=X_{(Y')}\ar[l]_{v'}\ar[d]^{u'}\\
    Y &
    Y'\ar[l]_{v}.
  }
  \tag{1.4.15.1}
  \label{III.1.4.15.1}
\]
Entonces, para todo $\sh{O}_X$-módulo cuasi-coherente $\sh{F}$, $\RR^q u_*'(\sh{F}')$ es canónicamente isomorfo a $\RR^q u_*(\sh{F})\otimes_{\sh{O}_Y}\sh{O}_{Y'}=v^*(\RR^q u_*(\sh{F}))$ para todo $q\geq 0$, donde $\sh{F}'={v'}^*(\sh{F})=\sh{F}\otimes_{\sh{O}_Y}\sh{O}_{Y'}$.
\end{proposition}

\oldpage[III]{93}
\begin{proof}
El homomorfismo canónico $\rho:\sh{F}\to v_*'({v'}^*(\sh{F}))$ (\textbf{0},~4.4.3.2) define por funtorialidad un homomorfismo
\[
  \RR^q u_*(\sh{F})\to\RR^q u_*(v_*'(\sh{F}')).
  \tag{1.4.15.2}\label{III.1.4.15.2}
\]

Por otra parte, poniendo $w=u\circ v'=v\circ u'$, se tienen los homomorfismos canónicos (\textbf{0}~,12.2.5.1 y 12.2.5.2)
\[
  \RR^q u_*(v_*'(\sh{F}'))\to\RR^q w_*(\sh{F}')\to v_*(\RR^q u_*'(\sh{F}')).
  \tag{1.4.15.3}\label{III.1.4.15.3}
\]

Componiendo (\hyperref[III.1.4.15.3]{1.4.15.3}) y (\hyperref[III.1.4.15.2]{1.4.15.2}), se obtiene un homomorfismo
\[
  \psi:\RR^q u_*(\sh{F})\to v_*(\RR^q u_*'(\sh{F}')),
\]
y se obtiene finalmente un homomorfismo canónico (cuya definición no impone \emph{hipótesis alguna} sobre $v$)
\[
  \psi^\sharp:v^*(\RR^q u_*(\sh{F}))\to\RR^q u_*'(\sh{F}'),
  \tag{1.4.15.4}
\]
y hay que probar que es un isomorfismo cuando $v$ es \emph{plano}.
Es claro que la cuestión es local sobre $Y$ y $Y'$; puede suponerse, por tanto, que $Y=\Spec(A)$ y $Y'=\Spec(B)$; se utilizará además el siguiente lema:
\begin{lemma}[1.4.15.5]
\label{III.1.4.15.5}
Sean $\vphi:A\to B$ un homomorfismo de anillos, $Y=\Spec(A)$, $X=\Spec(B)$, $f:X\to Y$ el morfismo correspondiente a $\vphi$, y $M$ un $B$-módulo.
Para que el $\sh{O}_X$-módulo $\widetilde{M}$ sea $f$-plano \sref[0]{0.6.7.1}, es necesario y suficiente que $M$ sea un $A$-módulo plano.
En particular, para que el morfismo $f$ sea plano, es necesario y suficiente que $B$ sea un $A$-módulo plano.
\end{lemma}

Esto se deduce de la definición \sref[0]{0.6.7.1} y de \sref[0]{0.6.3.3}, teniendo en cuenta \sref[I]{I.1.3.4}.

Así pues, de (\hyperref[III.1.4.11.1]{1.4.11.1}) y de las definiciones de los homomorfismos (\hyperref[III.1.4.15.3]{1.4.15.3}) (cf.~\sref[0]{0.12.2.5}) se deduce que $\psi$ corresponde entonces al morfismo compuesto
\[
  \HH^q(X,\sh{F})\xrightarrow{\rho_q}\HH^q(X,v_*'({v'}^*(\sh{F})))\xrightarrow{\theta_q}\HH^q(X',{v'}^*(v_*'({v'}^*(\sh{F}))))\xrightarrow{\sigma_q}\HH^q(X',{v'}^*(\sh{F})),
\]
donde $\rho_q$ y $\sigma_q$ son los homomorfismos de cohomología correspondientes a los morfismos canónicos $\rho$ y $\sigma:{v'}^*(v_*'(\sh{G}'))\to\sh{G}'$, y $\theta_q$ es el $\vphi$-morfismo (\textbf{0},~12.1.3.1) relativo al $\sh{O}_X$-módulo $v_*'({v'}^*(\sh{F}))$.
Pero, por la funtorialidad de $\theta_q$, se tiene el diagrama conmutativo
\[
  \xymatrix{
    \HH^q(X,\sh{F})\ar[rr]^{\rho_q}\ar[dd]_{\theta_q} & &
    \HH^q(X,v_*'({v'}^*(\sh{F})))\ar[dd]^{\theta_q}\\\\
    \HH^q(X',{v'}^*(\sh{F}))\ar[rr]^{{v'}^*(\rho_q)} & &
    \HH^q(X',{v'}^*(v_*'({v'}^*(\sh{F})))),
  }
\]
\oldpage[III]{94}
y, como por definición \sref[0]{0.4.4.3} ${v'}^*(\rho)$ es el inverso de $\sigma$, se ve que el morfismo compuesto considerado anteriormente no es finalmente otro que $\theta_q$; por consiguiente, $\psi^\sharp$ es el $B$-homomorfismo asociado $\HH^q(X,\sh{F})\otimes_A B\to\HH^q(X',\sh{F}')$.
Como $u$ es cuasicompacto, $X$ es una unión finita de abiertos afines $U_i$ ($1\leq i\leq r$); sea $\mathfrak{U}$ el recubrimiento $(U_i)$.
Como $v$ es un morfismo afín, también lo es $v'$ \sref[II]{II.1.6.2}[iii]; por consiguiente, los $U_i'={v'}^{-1}(U_i)$ forman un recubrimiento abierto afín $\mathfrak{U}'$ de $X'$.
Se sabe entonces \sref[0]{0.12.1.4.2} que el diagrama
\[
  \xymatrix{
    \HH^q(\mathfrak{U},\sh{F})\ar[r]^{\theta_q}\ar[d] &
    \HH^q(\mathfrak{U}',\sh{F}')\ar[d]\\
    \HH^q(X,\sh{F})\ar[r]^{\theta_q} &
    \HH^q(X',\sh{F}')
  }
\]
es conmutativo, y las flechas verticales son isomorfismos puesto que $X$ y $X'$ son esquemas \sref{III.1.4.1}.
Por consiguiente, basta probar que el $\vphi$-morfismo canónico $\theta_q:\HH^q(\mathfrak{U},\sh{F})\to\HH^q(\mathfrak{U}',\sh{F}')$ es tal que el $B$-homomorfismo asociado
\[
  \HH^q(\mathfrak{U},\sh{F})\otimes_A B\to\HH^q(\mathfrak{U}',\sh{F}')
\]
sea un isomorfismo.
Para toda sucesión $\mathbf{s}=(i_k)_{0\leq k\leq p}$ de $p+1$ índices de $[1,r]$, pónganse $U_\mathbf{s}=\bigcap_{k=0}^p U_{i_k}$, $U_\mathbf{s}'=\bigcap_{k=0}^p U_{i_k}'={v'}^{-1}(U_\mathbf{s})$, $M_\mathbf{s}=\Gamma(U_\mathbf{s},\sh{F})$ y $M_\mathbf{s}'=\Gamma(U_\mathbf{s}',\sh{F}')$.
La aplicación canónica $M_\mathbf{s}\otimes_A B\to M_\mathbf{s}'$ es un isomorfismo \sref[I]{I.1.6.5}; por tanto, la aplicación canónica $C^p(\mathfrak{U},\sh{F})\otimes_A B\to C^p(\mathfrak{U}',\sh{F}')$ es un isomorfismo, mediante el cual $d\otimes 1$ se identifica con la aplicación coborde $C^p(\mathfrak{U}',\sh{F}')\to C^{p+1}(\mathfrak{U}',\sh{F}')$.
Como $B$ es un $A$-módulo \emph{plano}, de la definición de los módulos de cohomología se deduce que la aplicación canónica $\HH^q(\mathfrak{U},\sh{F})\otimes_A B\to\HH^q(\mathfrak{U}',\sh{F}')$ es un isomorfismo \sref[0]{0.6.1.1}.
Este resultado se generalizará más adelante en \EGAUnavailableHyperref{section:III.6}{\textsection6}.
\end{proof}

\begin{corollary}[1.4.16]
\label{III.1.4.16}
Sean $A$ un anillo, $X$ un $A$-esquema de tipo finito y $B$ una $A$-álgebra fielmente plana sobre $A$.
Para que $X$ sea afín, es necesario y suficiente que $X\otimes_A B$ lo sea.
\end{corollary}

\begin{proof}
La condición es evidentemente necesaria \sref[I]{I.3.2.2}; se probará que es suficiente.
Como $X$ es separado sobre $A$ y el morfismo $\Spec(B)\to\Spec(A)$ es plano, de la Proposición~\sref{III.1.4.15} se deduce que
\[
\label{III.1.4.16.1}
  \HH^i(X\otimes_A B,\sh{F}\otimes_A B)=\HH^i(X,\sh{F})\otimes_A B
  \tag{1.4.16.1}
\]
para todo $i\geq 0$ y todo $\sh{O}_X$-módulo cuasi-coherente $\sh{F}$.
Si $X\otimes_A B$ es afín, el miembro izquierdo de (\hyperref[III.1.4.16.1]{1.4.16.1}) es nulo para $i=1$; por tanto, también lo es $\HH^1(X,\sh{F})$, puesto que $B$ es un $A$-módulo fielmente plano.
Como $X$ es un esquema cuasicompacto, la demostración termina por el criterio de Serre \sref[II]{II.5.2.1}.
\end{proof}

\begin{proposition}[1.4.17]
\label{III.1.4.17}
Sea $X$ un preesquema y sea $0\to\sh{F}\xrightarrow{u}\sh{G}\xrightarrow{v}\sh{H}\to 0$ una sucesión exacta de $\sh{O}_X$-módulos.
Si $\sh{F}$ y $\sh{H}$ son cuasi-coherentes, también lo es $\sh{G}$.
\end{proposition}

\oldpage[III]{95}
\begin{proof}
La cuestión es local sobre $X$, de modo que puede suponerse que $X=\Spec(A)$ es afín; basta entonces probar que $\sh{G}$ satisface las condiciones (d1) y (d2) de \sref[I]{I.1.4.1} (con $V=X$).
La verificación de (d2) es inmediata, pues, si $t\in\Gamma(X,\sh{G})$ se anula al restringirlo a $D(f)$, lo mismo ocurre con su imagen $v(t)\in\Gamma(X,\sh{H})$; por tanto, existe un $m>0$ tal que $f^m v(t)=v(f^m t)=0$ \sref[I]{I.1.4.1}, y, como $\Gamma$ es exacto por la izquierda, $f^m t=u(s)$, donde $s\in\Gamma(X,\sh{F})$; como $u$ es inyectivo, la restricción de $s$ a $D(f)$ es nula, luego por \sref[I]{I.1.4.1} existe un entero $n>0$ tal que $f^n s=0$; se deduce finalmente que $f^{m+n}t=u(f^n s)=0$.

Compruébese ahora (d1); sea $t'\in\Gamma(D(f),\sh{G})$; como $\sh{H}$ es cuasi-coherente, existe un entero $m$ tal que $f^m v(t')=v(f^m t')$ se prolonga a una sección $z\in\Gamma(X,\sh{H})$ \sref[I]{I.1.4.1}.
Pero, en virtud del Teorema~\sref{III.1.3.1} (o de (\textbf{I},~5.1.9.2)) aplicado al $\sh{O}_X$-módulo cuasi-coherente $\sh{F}$, la sucesión $\Gamma(X,\sh{G})\to\Gamma(X,\sh{H})\to 0$ es exacta; existe, pues, $t\in\Gamma(X,\sh{G})$ tal que $z=v(t)$; se ve así que $v(f^m t'-t'')=0$, denotando por $t''$ la restricción de $t$ a $D(f)$; por tanto, $f^m t'-t''=u(s')$, donde $s'\in\Gamma(D(f),\sh{F})$.
Pero, como $\sh{F}$ es cuasi-coherente, existe un entero $n>0$ tal que $f^n s'$ se prolonga a una sección $s\in\Gamma(X,\sh{F})$; como $f^{m+n}t'-f^n t''=u(f^n s')$, se ve que $f^{m+n}t'$ es la restricción a $D(f)$ de una sección $f^n t+u(s)\in\Gamma(X,\sh{G})$, lo que termina la demostración.
\end{proof}
